%% notes-tex/database/database-expansion-hosts/sections/rule.tex
%% [[experiments.database.expansion-hosts]]
%% https://github.com/Mjvolk3/torchcell/tree/main/notes-tex/database/database-expansion-hosts/sections/rule.tex
%% SOURCE: hand-authored prose. The rule, the rank quantity and the pool guard are read from
%% experiments/database/scripts/build_candidate_datasets_table.py and
%% experiments/database/scripts/build_bacteria_candidate_datasets_table.py; the per-row
%% figures come from the triage report
%% /scratch/projects/torchcell-scratch/tmp/deep-research/reports/Pk and next dataset wave triage.md

\section{The selection rule\secstatus{todo}}\label{sec:rule}

\subsection{Sequence reconstruction is the gate, not a preference\secstatus{todo}}\label{sec:gate}

A phenotype is only trainable against a genotype that can be written down, so every row names a
\term{sequence basis}, meaning the route by which the total genomic content of one strain would
be reconstructed. The bases are named per host, because the reference differs per host and because
a \texttt{PP\_} locus tag, a JGI \texttt{gNNNN} gene model, and a PA14 locus tag are different
namespaces that must stay distinguishable rather than being merged into one string field.

Two of the bases this pass needs do not exist in either pool today. The yeast
\texttt{SeqBasis} literal holds 11 values and every one of them is S288C-relative or
yeast-panel-relative, so \texttt{isolate-WGS} transfers to a \org{P. kudriavzevii} isolate
panel unchanged while a designed edit in the SD108 background has no value to take. The
bacterial literal is KT2440- and K-12-relative in the same way, so a transposon insertion in
\org{Stutzerimonas stutzeri} RCH2 or a guide library in PAO1 needs its own basis rather than a
borrowed one. Table~\ref{tab:schema} lists both additions with the policy each carries.

What fails the gate is a strain whose genome is unknown even in principle. The three
acetic-acid evolved \org{P. kudriavzevii} lines, PkAC-7, PkAC-8 and PkAC-9, land there: they
are characterized phenotypically and transcriptionally across six stresses, and no
resequencing, variant table, or accession surfaced in either search pass, so the genotype is
undefined and the strains are usable only as named strains with stress phenotypes. Rows in
that state are dropped in Table~\ref{tab:excluded} rather than ranked low, because no amount
of phenotypic value substitutes for a missing genotype.

\subsection{The rank quantity is measurements, not records\secstatus{todo}}\label{sec:measurements}

Rows are ordered by \term{measurements}, meaning records times phenotype dimensionality,
descending, which is the quantity \texttt{build\_candidate\_datasets\_table.py} already sorts
on. \term{Records} is the number of genotype $\times$ environment $\times$ replicate rows a
dataset can yield, and \term{dimensionality} counts the values each row carries. A scalar MIC
panel of $n$ isolates against $k$ drugs contributes $nk$ numbers, and a transcriptome panel of
$n$ isolates at 5,564 genes contributes $5.6\times10^{3}n$.

Ranking on records alone gets the order visibly wrong here, and the two clearest cases are
worth naming rather than asserting the rule. The 414-isolate clinical panel of Khaledi et al.
2020 holds 414 records, which would rank it below a 5,459-mutant arrayed screen, and yet each
of those records is a genome plus a transcriptome, so the row is roughly $2.4\times10^{6}$
measurements, inferred from the panel shape. The nine new pseudomonad rows hold 1,479 fitness
experiments, measured, against roughly $6.7\times10^{6}$ measurements, inferred, because each
experiment is a per-gene vector. The product places both correctly, and the supported-datasets
table applies that same normalization through its compressed-signal column.

\subsection{Band comes before tier, and why\secstatus{todo}}\label{sec:band}

Scale is the right default ordering and it is the wrong first question when a specific host is
being brought up. Each row therefore carries a \term{band} applied before tier and before
scale, and the composition by band is Table~\ref{tab:counts}.

The reason is measured in this pass rather than argued. The three regulatory-DNA rows at
$10^{8}$, $5.1\times10^{7}$ and $6.7\times10^{6}$ measured promoter sequences carry more
measurements than every other row combined, and none of the three bears on a low-data host,
because all three are \org{S. cerevisiae} and their perturbation is an 80 bp promoter insert in
a fixed scaffold. Conversely the largest \org{P. kudriavzevii} phenotype resource is 9,177
measurements, inferred, which is four orders of magnitude below the smallest pseudomonad
fitness row, so a pure measurement order puts every row of the new host below any plausible
cut. Both facts are consequences of what the hosts have been measured with rather than of what
they are worth to a transfer experiment, and a band corrects for exactly that difference.

\subsection{The per-host floor, and why a floor rather than a score bonus\secstatus{todo}}\label{sec:floor}

The recommended set carries a floor of rows per host, applied after ranking and reported row by
row, and every promotion it makes is printed in Table~\ref{tab:final} rather than folded into a
score.

A floor is the right instrument because the setup a new host needs is per host and is done
once. An \texttt{Organism} member, a host token in \file{torchcell/candidates/gates.py}, and a
reference release pinned in the genomes tier are each written the first time any row for that
host is built and never again, so the setup for a host with one built row is the same as the
setup for a host with twenty. A floor guarantees that every host in the wave starts, so that
one-time setup yields served records in this pass rather than in the next one.

A score bonus would do the same arithmetic and lose the audit. Folding a host weight into the
rank quantity makes a promoted row indistinguishable from a row that earned its position, and
the whole point of the measurement order is that a surprising position can be checked against
two numbers. The pin mechanism of the bacterial pass is the precedent: a stakeholder priority
is a real criterion, so it is applied and reported rather than absorbed into the score.

\subsection{Where the lines fall\secstatus{todo}}\label{sec:lines}

The long-run cut is 200 additions against the 170 datasets served, so the target is 370. The
pool ranked here does not reach the cut and the shortfall is reported rather than padded.
Sec.~\ref{sec:state-gaps} names the four axes the next discovery pass would search, and
Table~\ref{tab:asks} lists the nine asks that would convert unreachable rows into reachable
ones. Table~\ref{tab:counts} gives the composition of the pool by class and by host, so the
cost of the band order and of the floor is visible per host.

%% GENERATED FILE -- do not hand-edit.
%% SOURCE: experiments/database/scripts/build_hosts_candidate_datasets_table.py
\begin{table}[H]\centering
\small
\caption[]{Candidates by class and by host, split at the two wave lines.
\emph{Genotypes}, \emph{Instances} and \emph{Meas.} sum the per-row axes over the whole
pool; a row with no count contributes nothing, so every total is a lower bound.
\emph{Meas.} is instances times phenotype dimensionality, which is why the per-host
ordering of the measurement column does not follow the row counts: one yeast
regulatory-DNA row carries more measurements than every pseudomonad row together.}
\label{tab:counts}
\begin{tabular}{@{}l r r r r r r@{}}
\toprule
 & Wave 1 & Wave 2 & Reserve & Genotypes & Instances & Meas. \\
\midrule
CRISPR library screen & 1 & 0 & 6 & 96,055 & $1.3\times 10^{5}$ & $1.3\times 10^{5}$ \\
Deep mutational scan & 3 & 0 & 7 & 925,159 & $1.5\times 10^{6}$ & $1.5\times 10^{6}$ \\
Engineering parts & 3 & 1 & 1 & 3,116 & $3.1\times 10^{3}$ & $3.1\times 10^{3}$ \\
Essentiality & 0 & 0 & 7 & 7,385 & $2.3\times 10^{3}$ & $2.3\times 10^{3}$ \\
Fitness / chemical genomics & 0 & 9 & 10 & 2,753,324 & $2.5\times 10^{7}$ & $3.2\times 10^{7}$ \\
Metabolome / flux & 0 & 0 & 8 & 359 & $9.6\times 10^{3}$ & $1.9\times 10^{6}$ \\
Modality / backbone & 0 & 1 & 2 & 14,247 & $1.4\times 10^{4}$ & $1.4\times 10^{4}$ \\
Natural variation & 4 & 0 & 15 & 28,291 & $4.1\times 10^{5}$ & $4.2\times 10^{5}$ \\
Production campaign & 0 & 5 & 26 & 471 & $1.5\times 10^{3}$ & $1.7\times 10^{3}$ \\
Proteome & 1 & 1 & 13 & 4,028 & $5.8\times 10^{3}$ & $7.9\times 10^{5}$ \\
Regulatory DNA & 27 & 1 & 0 & 166,789,446 & $1.7\times 10^{8}$ & $2.1\times 10^{8}$ \\
Tolerance / robustness & 0 & 1 & 11 & 8,267 & $9.5\times 10^{3}$ & $2.0\times 10^{5}$ \\
Transcriptome & 7 & 1 & 18 & 477 & $3.4\times 10^{3}$ & $1.6\times 10^{7}$ \\
Transposon fitness & 4 & 0 & 7 & 105,459 & $3.4\times 10^{6}$ & $9.8\times 10^{6}$ \\
\midrule
\multicolumn{7}{@{}l}{\emph{The same rows, by host}}\\
E. coli & 9 & 0 & 16 & 2,916,904 & $6.0\times 10^{6}$ & $5.5\times 10^{7}$ \\
P. aeruginosa & 0 & 1 & 23 & 109,342 & $1.1\times 10^{5}$ & $2.6\times 10^{6}$ \\
P. fluorescens & 0 & 4 & 1 & 100,000 & 745 & $3.4\times 10^{6}$ \\
P. kudriavzevii & 8 & 10 & 21 & 391 & $9.9\times 10^{3}$ & $3.6\times 10^{5}$ \\
P. putida & 8 & 0 & 41 & 10,341 & $5.5\times 10^{4}$ & $5.7\times 10^{6}$ \\
P. simiae & 0 & 1 & 0 & 0 & 201 & $9.1\times 10^{5}$ \\
P. syringae & 0 & 2 & 0 & 0 & 69 & $3.1\times 10^{5}$ \\
S. cerevisiae & 25 & 1 & 29 & 167,599,106 & $1.9\times 10^{8}$ & $2.0\times 10^{8}$ \\
S. stutzeri & 0 & 1 & 0 & 0 & 350 & $1.6\times 10^{6}$ \\
\midrule
\multicolumn{7}{@{}l}{\emph{The same rows, by band}}\\
regulatory DNA & 34 & 1 & 0 & 166,797,405 & $1.7\times 10^{8}$ & $2.1\times 10^{8}$ \\
cross-host fitness & 0 & 9 & 0 & 0 & $1.4\times 10^{3}$ & $6.2\times 10^{6}$ \\
host bootstrap & 8 & 10 & 21 & 391 & $9.9\times 10^{3}$ & $3.6\times 10^{5}$ \\
scale & 8 & 0 & 110 & 3,938,288 & $3.0\times 10^{7}$ & $5.5\times 10^{7}$ \\
\midrule
Total & 50 & 20 & 131 & 170,736,084 & $2.0\times 10^{8}$ & $2.7\times 10^{8}$ \\
\bottomrule
\end{tabular}
\end{table}


One line is set by the machinery rather than by judgment. The yeast pool holds 169
\texttt{Candidate} literals, of which 154 are unbuilt and ranked, and its \texttt{main()}
raises on a pool over 250 rows, so a 200-row wave cannot go into one pool at all. The guard
forces the split by organism, and the gate needs the same split for its own reasons. The yeast
\texttt{Candidate} model has no \texttt{organism} field at all, the bacterial one admits two
values, and \texttt{verdict.CandidateTable} admits only the two strings
\texttt{bacteria} and \texttt{yeast}, which is the \texttt{table} field of all 52 stored
verdicts. Table~\ref{tab:schema} records each of those edits.
